\documentclass[11,5pt, a4paper]{article}
\usepackage[a4paper, total={6in, 9.5in}]{geometry}
\usepackage{float}
\usepackage{graphicx} % Required for inserting images
\usepackage{natbib} % bibliografia
\usepackage{url}
\usepackage{booktabs} % tabella
\usepackage{amsmath}

\usepackage{titling}
\usepackage[italian]{babel}
\usepackage[utf8]{inputenc} % Per caratteri accentati
\usepackage[T1]{fontenc}    % Per font corretti in output

\setlength{\droptitle}{-3em}

\usepackage{etoolbox}% http://ctan.org/pkg/etoolbox
\makeatletter
\let\oldtableofcontents\tableofcontents
\renewcommand{\tableofcontents}{\begingroup%
  \patchcmd{\@makeschapterhead}% <cmd>
    {\vspace*{50\p@}}% <search>
    {}% <replace>
    {}{}% <success><failure>
  \oldtableofcontents%
  \endgroup%
}
\makeatother

\title{\normalsize Introduzione alla Data Science e al Pensiero Computazionale \\[0.5em]
\normalsize Gruppo 14 - Arianna Bertoldo e Veronica Contarin \\[1.3 em]
\LARGE \textbf{Data-driven marketing: analisi di dati bancari \\ per la sottoscrizione di un conto deposito}}

\date{}
\author{}

\begin{document} 

\maketitle

\vspace{-5.6em}
\tableofcontents %COME SI METTE IN ITALIANO :(

\section{Introduzione}
% che attacco tremendo
% il telemarketing spacca
Il telemarketing è una pratica diffusa in numerosi settori, tra cui quello bancario e finanziario. Consiste nel contattare direttamente potenziali clienti per promuovere prodotti o servizi. 
Sebbene sia un metodo rapido e potenzialmente efficace, comporta costi organizzativi ed economici significativi, con il rischio di disperdere risorse su contatti infruttuosi. La sua efficacia dipende infatti dalla capacità di individuare a monte i clienti probabilmente più propensi all'adesione o all'acquisto e contattarli nel momento più adeguato.

A questo proposito, la crescente disponibilità di dati consente e incoraggia le aziende a ricorrere a un approccio \textit{data-driven} per ottimizzare le loro attività.
Le informazioni raccolte nelle campagna possono orientare l'organizzazione delle campagne successive, selezionando accuratamente i destinatari e personalizzando le comunicazioni. L'analisi dei comportamenti e delle caratteristiche della clientela può così contribuire a ridurre i contatti infruttuosi e a rendere il telemarketing più mirato ed efficiente \citep{ROSARIO2023}.
\\

In questa relazione, presentiamo un'analisi dei dati \citep{ariannab02}
che abbiamo condotto sulla base dei risultati di una campagna di telemarketing bancaria mirata alla sottoscrizione di conti deposito vincolati.  L'obiettivo ultimo era quello di predire quali clienti avrebbero sottoscritto il prodotto proposto, costruendo modelli predittivi per identificare i clienti che presentano una maggiore probabilità di sottoscrivere il conto deposito. 


\section{Dataset}

Il dataset esaminato è \textit{Bank\_Marketing.csv} \citep{BankMarketing}.
Si tratta di un insieme di dati raccolti da una banca portoghese tra maggio 2008 e novembre 2010, nel corso di una campagna marketing per la sottoscrizione di un conto deposito vincolato \citep{bank_marketing_222}.

Il dataset si compone di 45211 osservazioni, per le quali sono raccolte informazioni relative a 17 variabili, di cui 16 variabili esplicative e una variabile target (\textit{y}). 
Le variabili esplicative possono essere suddivise in quattro gruppi. Il primo comprende le caratteristiche personali e socio-economiche del cliente, quali età, occupazione, stato civile e livello di istruzione. Il secondo riguarda la sua situazione finanziaria, attraverso il saldo medio annuo, la presenza di prestiti personali o immobiliari e l'eventuale inadempienza creditizia. Il terzo gruppo descrive la campagna corrente, includendo il mezzo di contatto, il giorno e il mese dell'ultima chiamata, la sua durata e il numero complessivo di contatti effettuati. Infine, alcune variabili riportano informazioni sulla precedente campagna, come il numero di contatti, il tempo trascorso dall'ultima chiamata e il relativo esito.



\begin{table}[h]
\centering
\caption{Classificazione delle variabili per area tematica}
\begin{tabular}{@{}ll@{}}
\toprule
\textbf{Gruppo} & \textbf{Variabili} \\ \midrule
Socio-demografiche & age, job, marital, education \\
Finanziarie & balance, default, loan, housing \\
Campagna corrente & contact, day, month, duration, campaign \\
Campagna precedente & pdays, previous, poutcome \\
Target & y \\ \bottomrule
\end{tabular}
\end{table}

Il dataset è ben strutturato, ma in alcune categorie presenta numerosi valori mancanti, categorizzati come \textit{unknown}. Per la variabile \textit{contact}, il 29\% dei valori sono assenti, poiché la raccolta di tale dato è iniziata dopo due mesi dall'avvio della campagna. Ulteriormente, anche \textit{job} e \textit{education} presentano delle basse quote (inferiori al 5\%) di valori mancanti. Infine, per \textit{poutcome}, quasi l'82\% delle osservazioni sono identificate come sconosciute; tuttavia, ciò non è un segnale di bassa qualità dei dati, bensì è legato alla natura della variabile: le osservazioni mancanti sono per lo più riconducibili a clienti che non sono stati contattati nella precedente campagna, e di conseguenza l'esito non è identificabile.

In seguito è stato notato che, in mancanza di un identificativo univoco per i singoli clienti contattati, c'è la possibilità che talune persone siano state incluse nel dataset più volte, a seguito di ripetute chiamate. In questo caso, le osservazioni non sono mai esattamente corrispondenti, se non per le caratteristiche socio-demografiche e finanziarie.

\section{Metodologia}
Il lavoro di analisi dei dati è stato organizzato in tre fasi. Nella prima è stata realizzata un'analisi descrittiva del dataset, mirata a comprenderne la struttura, la qualità e le distribuzioni delle variabili. Nella seconda fase sono state esaminate le relazioni tra le variabili esplicative e la variabile target, con il supporto di visualizzazioni grafiche. Infine, nella terza fase, i dati sono stati preparati per la classificazione e sono stati confrontati diversi algoritmi di apprendimento supervisionato.

\subsection{Fase 1: Analisi descrittiva ed esplorativa}
\label{ss: Fase 1}
La prima fase ha riguardato l'analisi della struttura del dataset, dei tipi di variabile, della presenza di informazioni mancanti, e di alcune relazioni tra le variabili. 

Per le variabili numeriche sono state calcolate statistiche descrittive, come misure di centralità e di dispersione. Per le variabili categoriche sono state invece considerate le frequenze assolute e relative.

Tramite l'analisi, sono emersi dei potenziali limiti del dataset, legati prevalentemente a valori mancanti (\textit{unknown}) e allo sbilanciamento della variabile target.

Ulteriormente, nella prima fase, sono state avanzate domande e ipotesi relative alle caratteristiche legate alla sottoscrizione. Queste relazioni sono state esplorate tramite confronti tra categorie e gruppi di clienti, e coefficienti di correlazione. In particolare, sono state esaminate le differenze nella propensione alla sottoscrizione secondo caratteristiche socio-demografiche e finanziarie, al numero di contatti effettuati, all'esito della precedente campagna e alla durata delle chiamate.

\subsection{Fase 2: Analisi delle relazioni tra le variabili}
\label{ss:Fase2}
Nella seconda fase, l'analisi si è concentrata sulle relazioni tra le variabili esplicative e la variabile target, ossia l'esito della campagna. 
L'obiettivo era individuare quali caratteristiche dei clienti e quali aspetti dei contatti fossero associati a un esito positivo della campagna. Le analisi svolte in questa sezione si sono innestate sui risultati e sulle riflessioni della Fase 1.

Sono stati confrontati i tassi di adesione tra diverse categorie lavorative, situazioni finanziarie, fasce d'età e mezzi di contatto. Inoltre, sono state esaminate le relazioni tra l'esito della campagna e le variabili relative alla precedente campagna, tra cui il numero di chiamate già effettuate, il numero di contatti ed esito della campagna precedente. Un'ulteriore aspetto esplorato è stata la durata delle chiamate, con la sua distribuzione e andamento medio, comparandole in base all'esito della campagna. 
A questo proposito, per porre enfasi sulla successione temporale, la variabile \textit{month} è stata trasformata in un indice progressivo. In tal modo, mesi con lo stesso nome ma appartenenti ad anni diversi non sono stati ridotti alla stessa categoria, così facilitando la rappresentazione dell'evoluzione dell'organizzazione della campagna nel tempo.

Infine, si osserva che le categorie più positivamente associate alla variabile target sono sottorappresentate nel dataset, limitando la portata dei risultati.

\subsection{Fase 3: Modelli di classificazione}
\label{ss: Fase 3}
In questa fase si sono confrontati tre modelli di classificazione supervisionata: regressione logistica, k-Nearest Neighbours e Random Forest.

Il primo modello, la regressione logistica, stima la probabilit\`{a} che un'osservazione appartenga alla classe positiva, esaminando la relazione tra variabili predittive e variabile target \citep{REGLOG_IBM}. Per normalizzare le variabili sulle quali \`{e} stato allenato il modello, \`{e} stata utilizzata la seguente formula per la standardizzazione:

\begin{equation*}
    Z = \frac{X-\mu}{\sigma}
\end{equation*}
\vspace{0.22em}

Il secondo modello realizzato è k-Nearest Neighbours. Questo modello assegna una classe a ciascuna osservazione in base alle classi prevalenti tra le $k$ osservazioni più vicine all'interno del dataset. Dato che la sua performance dipende dal numero di vicini presi in considerazione, sono stati confrontati vari valori di $k$ e per ciascuno di questi sono stati valutati gli effetti su f1-score per la classe positiva e accuracy complessiva. Sulla base dei risultati ottenuti, il modello finale è stato addestrato utilizzando $k=3$.

Infine, è stato considerato il modello Random Forest. Tale modello è composto da un insieme di alberi decisionali, ciascuno addestrato su differenti campioni e sottoinsiemi di variabili. La combinazione di più alberi rende possibile arrivare a previsioni più robuste rispetto a quelle prodotte da un singolo albero decisionale, grazie all'ensemble learning \citep{RF_IBM}.
Per la Random Forest sono state confrontate diverse quantità di stimatori e soglie di classificazione. Alla fine, si è optato per un modello con 100 alberi e una soglia di 0.6. Inoltre, sono stati aggiustati alcuni parametri, come la profondità, e si è conferito un peso bilanciato alle osservazioni (\textit{weight=balanced\_subsample}), in modo da assegnare un peso maggiore alla classe positiva che è sottorappresentata nel dataset.

\section{Dati e interpretazione}

\subsection{Fase 1 e 2: Domande di ricerca e analisi dati}
Dalle analisi esplorative è emerso fin da subito come la classe target (i clienti che hanno aperto un conto deposito presso la banca) sia in forte minoranza rispetto ai clienti non aderenti all'interno del campione: al termine della raccolta dei dati, solo l’11\% del totale ha sottoscritto un conto deposito. 

\subsubsection{Classi lavorative, giacenza media, livello di istruzione}

Ci si è domandati se l'appartenenza a determinate classi lavorative potesse portare ad un maggiore desiderio di sottoscrivere ad un conto deposito vincolato.
È emerso che gli studenti e i pensionati sono state le classi lavorative più aderenti, rispettivamente il 28,7\% e il 22,8\% di loro. Si è verificato che giovani e anziani avessero effettivamente tassi di iscrizione maggiori rispetto alla fascia intermedia d'età, e così è risultato: il 33,4\% dei clienti over 60 aveva aperto un conto deposito, così anche il 17,6\% degli under 30, e solo il 10\% di chi aveva fra i 30 e i 59 anni. 
Tuttavia giovani e anziani costituiscono la minor parte del campione, rispettivamente l'11,6\% e il 4\%. 

\begin{figure}[H]
\centering
\includegraphics[width=0.5\textwidth]{GRAFICI PER LATEX/subscription-jobs.png}
\caption{Distribuzione del tasso di sottoscrizione a conto deposito per classi lavorative}
\label{fig:job}
\end{figure}

Gli operai sono stati il gruppo professionale meno aderente, con appena il 7,3\% di sottoscrizioni, seguiti da imprenditori, casalinghi e lavoratori nei servizi, tutti con un tasso di adesione che non supera il 9\%.
Ipotizzando che una minor propensione sia dovuta ad un fattore comune a queste classi lavorative, abbiamo vagliato due ipotesi principali: \textit{è la disponibilità economica ad agire sulla propensione alla sottoscrizione o il livello di istruzione?}

Nessuna delle due è risultata in grado di spiegare il comportamento della variabile target (tenuto conto, cioè, che i più propensi ad aprire un conto deposito sono stati studenti, pensionati, disoccupati, manager e amministrativi).
È vero che per le categorie con saldi medi annui più alti la campagna ha avuto più successo (pensionati), ma non si tratta di un rapporto causa-effetto. Gli studenti, ad esempio, sono una delle due classi più aderenti: pur essendo i più risparmiatori (i loro conti sono i meno in negativo fra tutte le classi lavorative), non sono i più ricchi. Le categorie con saldi medi annui più bassi sono effettivamente meno interessate ad aprire un conto deposito (vedi operai), ma non vi è un trend preciso e identificabile.

Non notiamo una particolare relazione nemmeno fra livello di istruzione e tasso di successo della campagna. I più istruiti con istruzione di livello universitario sono risultati manager, lavoratori autonomi e imprenditori, mentre lavoratori nei servizi e tecnici sono stati le classi con le percentuali più alte di diplomati. Tuttavia, questi gruppi non si sono distinti per tassi pià alti di sottoscrizione allo strumento finanziario propostogli dalla banca.
E in effetti i clienti con un livello di istruzione superiore alla scuola primaria risultano solo leggermente più propensi a sottoscrivere un conto deposito. In particolare, il 15\% di chi possiede un’istruzione universitaria o equiparabile è iscritto, contro l’8,6\% di chi possiede la sola licenza elementare.

\begin{figure}[H]
\centering
\includegraphics[width=0.6\textwidth]{GRAFICI PER LATEX/education-jobs.png}
\caption{Distribuzione del livello di istruzione per classi lavorative}
\label{fig:education-job}
\end{figure}

Concludiamo, allora, che studenti e pensionati siano i profili più in linea con lo strumento finanziario del conto deposito vincolato. Ad iscriversi potrebbe spingerli la ricerca di uno strumento semplice e affidabile per i propri risparmi, ma questo aspetto non è indagabile da una ricerca come quella svolta in questa sede.

\subsubsection{Campagna precedente}
Altri aspetti sono risultati più rilevanti. È stato studiato il rapporto tra la campagna precedente e quella corrente.
I clienti che già hanno sottoscritto un conto deposito hanno tassi di successo molto più alti rispetto agli altri gruppi.
Il 23\% dei clienti già contattati nella campagna marketing precedente hanno aderito all'offerta proposta, contro il 9\% dei clienti mai contattati.
Se la scorsa campagna aveva già avuto successo, la percentuale di sottoiscrizioni in quella attuale arriva al 64,7\%. I clienti aderenti ma non convinti nella campagna passata sono invece il 12,6\%.
Inoltre, un alto numero di contatti nella storia del cliente è associato ad un tasso di iscrizione più alto.

\begin{figure}[H]
\centering
\includegraphics[width=0.6\textwidth]{GRAFICI PER LATEX/contact-frequency-poutcome-success.png}
\caption{Tassi di successo e insuccesso per numero di contatti ed esito della campagna precedente}
\label{fig:frequency-contact-poutcome}
\end{figure}

Clienti già contattati in passato, a maggior ragione se hanno già aderito ad un'offerta della banca, dimostrano una fiducia nei confronti della istituto finanziario che si concretizza in tassi di adesione più alti alle sue campagne marketing.

Diversamente da quella precedente, nella campagna attuale, la probabilità di successo è massima per i clienti contattati una singola volta: il 14,6\% per una chiamata fino a scemare al 4\% per più di undici chiamate.
La durata media di una chiamata è inferiore ai cinque minuti. La durata massima è stata di un'ora e venti.
Tolti i casi limite, le chiamate più lunghe sembrano aggirarsi attorno agli undici minuti e quelle standard sui tre.

\begin{figure}[H]
\centering
\includegraphics[width=1\textwidth]{GRAFICI PER LATEX/duration-success-rate.png}
\caption{Andamento del tasso successo e di insuccesso delle chiamate per durata}
\label{fig:duration-success}
\end{figure}

Emerge una chiara associazione tra durata delle chiamate e tasso di successo: le chiamate con tassi di adesione più alti sono quelle con durate nettamente maggiori. Questo andamento è osservabile soprattutto nei primi tre quarti della campagna, ma persiste anche nella fase finale, dove l'andamento delle chiamate di successo si conforma a quello delle chiamate senza successo. 
Ipotizziamo che a quel punto il numero di clienti disponibili ad essere chiamati sia diminuito ulteriormente, portando ad un allineamento con la durata media delle chiamate per il campione (campione, ricordiamo, fortemente sbilanciato per la sovrabbondanza di clienti non aderenti).

Abbiamo notato infatti che man mano che le chiamate si allungano queste sono associate a percentuali di sottoscrizioni più alte: a cinque minuti, almeno un cliente su dieci sottoscrive; a dieci minuti, sono diventati quattro;  ai quindici, sei clienti su dieci risultano aderenti.

\subsubsection{Periodo della campagna}

Il maggior numero di chiamate è concentrato nel primo anno, con picchi intorno ad alcuni intervalli. La media in questo primo periodo è di circa 3000 chiamate al mese. Nel secondo anno la media delle chiamate effettuate ogni mese si abbassa alle 300. Ottobre del primo anno è il mese col maggior tasso di successo (61\%), con circa 1500 chiamate effettuate.
Gli altri mesi di successo sono tutti concentrati nel secondo anno, come marzo del “terzo anno”, ma ipotizziamo che questo sia dovuto al minor numero di chiamate effettuate a quel punto.

\begin{figure}[H]
\centering
\includegraphics[width=0.6\textwidth]{GRAFICI PER LATEX/monthly-calls.png}
\caption{Distribuzione chiamate effettuate mensilmente}
\label{fig:monthly-calls}
\end{figure}

\subsubsection{Correlazione e Feature Selection}

Dallo studio dei coefficienti di correlazione, emerge effettivamente una relazione mediamente forte di segno positivo fra \textit{duration} e la variabile target. 

\begin{table}[H]
    \centering
    \begin{tabular}{|l|r|r|}
    \hline
    Metodo    & Pearson & Spearman\\
    \hline
    duration  & 0.395 & 0.342\\
    \hline
    pdays     & 0.104 & 0.154\\
    \hline
    previous  & 0.093 & 0.169\\
    \hline
    balance   & 0.053 & 0.100\\
    \hline
    age       & 0.025 & -0.009\\
    \hline
    day       & -0.028 & -0.030\\
    \hline
    campaign  & -0.073 & -0.084\\
         \hline
    \end{tabular}
    \caption{Coefficienti di correlazione con \textit{y}}
    \label{tab:correlation-table}
\end{table}

Dallo studio dei coefficienti di correlazione, emerge effettivamente una relazione mediamente forte di segno positivo fra \textit{duration} e la variabile target. Questa relazione è la più marcata rilevata nel dataset. Possiamo ipotizzare che chiamate più lunghe equivalgano ad un dialogo prolungato sulle condizioni e le caratteristiche dell’offerta, oltre che alla finalizzazione del contratto.
È fondamentale sottolineare che questa non sia una variabile propriamente predittiva, poiché è un dato che viene registrato solo in seguito alla comunicazione col cliente e non è conoscibile in anticipo. 

Notiamo differenze per le due tipologie di indici: Pearson individua valori positivi, anche se meno forti, con la frequenza dei contatti e l’essere già stati contattati.
Sapendo che i valori dei coefficienti di correlazione variano tra -1 e 1, i valori inferiori a quelli menzionati sono trascurabili.
Spearman, d’altra parte, cattura valori leggermente più alti per quasi tutte le feature, con l’eccezione di \textit{duration} (-0,052) e di \textit{age} (-0,033).
In questo caso il numero di contatti nella campagna precedente per singolo cliente è più influente (+0,075).
La correlazione fra saldo in banca e variabile target raddoppia, pur rimanendo piuttosto bassa (0,1).
Questi sono gli effetti della robustezza di Spearman agli outlier, evidentemente presenti nelle variabili considerate.

Dal momento che i coefficienti erano stati calcolati prima di creare la variabile \textit{month\_index}, prendiamo in considerazione anche i risultati della \textit{feature importance} rilevata dal modello Random Forest.
 In questo caso possiamo notare che sia la durata della chiamata e sia il mese in cui è stata effettuata sono risultate le variabili più discriminanti nelle scelte degli alberi. 
Tuttavia, si tratta di fattori poco sfruttabili dal punto di vista predittivo. Da un lato la durata delle chiamate che non è conoscibile a priori; dall'altro il mese in cui la chiamata viene effettuata, la cui influenza sul valore della variabile target è interpretabile solo nel contesto della campagna marketing svolta dalla banca portoghese.

\subsection{Fase 3: Confronto tra modelli}
Utilizzando tre diversi modelli di classificazione binaria e confrontando gli indicatori di performance, osserviamo che il forte sbilanciamento delle classi, per cui la classe 0 è nettamente più numerosa della classe 1, ha inficiato le capacità di predizione di tutti i modelli (\textit{support} per la classe 1 è di appena duemila casi, contro i quasi sedicimila della classe 0). Difatti tutti i modelli hanno performance considerevolmente peggiori per la classe 1, ovvero i potenziali clienti aderenti.


\begin{table}[H] 
    \centering
    \begin{tabular}{|l|l|l|l|l|l|}
    \hline
    Modello predittivo     & Accuracy & Precision & Recall & F1-score & Support\\
    \hline
    Regressione logistica     & 0.90628 & 0.66746 & 0.39650 & 0.49748 & 2116\\
    \hline
    K-Nearest Neighbours    & 0.89113 & 0.55261 & 0.36484 & 0.43951 & 2116\\
    \hline
    Random Forest     & 0.89964 & 0.55335 & 0.73771 & 0.63237 & 2116\\
    \hline
    \end{tabular}
    \caption{Confronto performance modelli, valori classe 1}
    \label{tab:models}
\end{table}

\begin{itemize}
    \item l’\textit{accuratezza} è sostanzialmente equivalente per Regressione Logistica e Random Forest: entrambi i modelli sono stati in grado di predire correttamente nove casi su dieci nel dataset di test (ma la maggioranza di questi casi sono clienti della classe 0, riferita a chi non ha sottoscritto il conto deposito)
\end{itemize}

\begin{itemize}
    \item la Regressione Logistica è il modello più \textit{preciso} nell’individuare i clienti in classe 1: identifica il 66\% dei casi positivi
\end{itemize}

\begin{itemize}
    \item per la \textit{recall} per la classe 1 il modello con risultati migliori è Random Forest: ha identificato il 73\% dei clienti effettivamente iscritti, contro il 36\% degli altri due modelli
\end{itemize}

\begin{itemize}
    \item è sempre Random Forest ad ottenere il valore di \textit{f1-score} più alto
\end{itemize}

Concludiamo che, per il nostro dataset, Random Forest è il modello più adatto a prevedere i potenziali clienti; riteniamo che il risultato ottenuto per f1-score sia un compromesso accettabile. Resta il fatto che il modello non eccelle, ma è capace di gestire un dataset fortemente sbilanciato a favore dei clienti non iscritti, dimostrando un'ottima capacità predittiva per i clienti realmente aderenti.


\section{Conclusioni}
 
Nel complesso, l'analisi del dataset \textit{Bank\_Marketing.csv} ha permesso di individuare gli elementi maggiormente associati alla sottoscrizione del conto deposito. Tra questi, la durata della chiamata emerge come il principale indicatore di un forte interesse da parte del cliente, affiancato dall'esito positivo della campagna precedente, che segnala fiducia e stima verso l'istituto bancario.

A questi elementi si aggiungono fattori socio-demografici e finanziari rilevanti, come la maggiore propensione alla sottoscrizione riscontrata tra studenti e over 60, profili che sembrano favorire strumenti di risparmio semplici e affidabili (ma che sono presenti in percentuali minime nel dataset). Si è inoltre riscontrata una lieve tendenza positiva associata all'assenza di debiti personali o mutui. Il periodo in cui avviene il contatto è risultato anch'esso molto rilevante, rivelandosi la seconda variabile più importante nel modello Random Forest; tuttavia, come ipotizzato, la sua influenza è più che altro legata alla gestione organizzativa della campagna, con contatti distribuiti in modo disomogeneo nell’arco di tempo considerato, piuttosto che a caratteristiche dei soggetti contattati. Per tale ragione, pur essendo un dato utile a descrivere l'andamento passato, questo fattore risulta possibilmente meno sfruttabile per la previsione di scenari futuri.

Questi risultati indicano che il responso del cliente non dipenda solo da caratteristiche personali, ma che sia invece influenzato anche dalla storia delle sue precedenti relazioni con la banca, e dal modo in cui la campagna viene gestita.

I modelli predittivi analizzati hanno ottenuto risultati piuttosto simili tra loro, con prestazioni discrete. La Random Forest ha mostrato i punteggi più incoraggianti, specialmente nella capacità di individuare i potenziali sottoscrittori (recall del 73\%), pur arrancando per la precisione. \\

Per testare la correttezza delle nostre conclusioni, abbiamo voluto confrontare quanto detto finora con altri studi svolti sul tema. Ciò è stato possibile grazie alla folta letteratura di studi sui modelli predittivi applicati al telemarketing bancario. Fra questi molti erano basati sul medesimo dataset da noi utilizzato (con il fine di predire l’adesione ad un conto deposito), con alcune differenze tra caso e caso. In particolare, alcuni dataset presentavano feature aggiuntive (peraltro emerse come influenti nella predizione della variabile target) quali alcuni indicatori macroeconomici: fra questi Euribor (un tasso di riferimento giornaliero sui tassi di interesse delle banche nella zona euro), la variazione del tasso di occupazione e l'indice dei prezzi al consumo.

Un’altra differenza rilevante sono state le tecniche di machine learning implementate. \citet{Zhang_RF_2022} per esempio, introducono un framework ibrido che utilizza l’\textit{ensemble modeling} (combinazione di più modelli) per una gestione ottimale di dataset sbilanciati, ottenendo risultati migliori rispetto agli approcci tradizionali con un singolo modello. Difatti, il modello che ha dimostrato performance migliori sul nostro dataset è stato Random Forest, che utilizza più alberi decisionali in combinazione.
Pare, infatti, che lo sbilanciamento fra classi sia una caratteristica intrinseca dei dataset sui profili bancari, per cui le sottoscrizioni sono sempre una minoranza sul totale. In aggiunta a questa difficoltà, le variabili in gioco sono molte e differenziate, ponendo un ulteriore livello di complessità alla sfida di predire correttamente il successo della campagna \citep{peter_2025}. \\

Uno studio \citep{medium_2025} molto simile al nostro (basato sul dataset UCI \citep{bank_marketing_222}) ci ha permesso di confrontare i risultati da noi ottenuti e verificarne la veridicità.
Vengono così confermate le nostre osservazioni relative all’impatto della durata delle chiamata (per cui chiamate più lunghe sono associate a livelli più alti di sottoscrizione) e alla frequenza dei contatti (per cui minori contatti in questa campagna sono associati a tassi di adesioni più alti).
Aver contattato ripetutamente il cliente non ha aumentato le probabilità di successo della campagna, ma l’aver già avuto comunicazioni con un cliente nella campagna precedente sì. I modelli utilizzati sono stati i medesimi, ma le tecniche di scaling diverse (Robust e MinMax). In questo studio il punteggio 'di compromesso' per precisione e recall si attesta attorno al 69\%, contro il nostro f1-score del 63\% ottenuto con Random Forest (recall al 73\% e precision al 55\%, che si è attestato il modello migliore nella nostra analisi. \\

Sebbene i risultati non siano ideali, vanno contestualizzati: riteniamo che per una banca, soprattutto se di grandi dimensioni e con fondi maggiori da riservare alla ricerca e al marketing strategico, sia preferibile fare molte chiamate a vuoto, piuttosto che perdere potenziali clienti interessati. Infatti, una precision bassa equivale a tanti falsi positivi, cioè clienti contattati ma erroneamente individuati dal modello come probabili aderenti. Una recall più alta, invece, è equivalente ad una quota minore di falsi negativi, cioè clienti mai contattati perché non selezionati dal modello come potenziali sottoscriventi. \\


In definitiva, la nostra analisi potrebbe essere utile per informare decisioni in merito al targeting di una campagna marketing come quella oggetto del dataset. A questo proposito, i modelli predittivi realizzati non sono sufficienti a prevedere con elevata probabilità l'esito della campagna se usati da soli. Ciononostante, possono costituire un utile supporto alle decisioni riguardo a quali categorie dare la priorità e a come impostare la campagna.


\section{Bibliografia}
\bibliography{biblio}
\bibliographystyle{plainnat}

\end{document}
